\section{Metodologia}
\label{sec:metodologia}
\subsection{Seleção e coleta}
A população de candidatos é obtida pela Search API do GitHub, com pelo menos
mil estrelas, sem forks nem repositórios arquivados e com atividade desde o
início da janela. A janela é o intervalo UTC de 1 de outubro de 2025 (inclusivo)
a 1 de outubro de 2026 (exclusivo). Incluímos projetos com workflows próprios
em \texttt{.github/workflows/}, pelo menos cinco releases estáveis publicadas e cinquenta
runs válidos dentro da janela. Pré-releases são preservadas para variantes, mas ficam fora da definição
principal de deploy e inclusão; drafts são excluídos.

A coleta registra metadados, releases e workflow runs do default branch,
com evento \texttt{push}. As consultas de runs são mensais; a integração
subdivide intervalos saturados para lidar com o teto de mil resultados por
consulta. Os IDs são deduplicados. Cache por página e checkpoints permitem
retomada, e a incompletude é registrada, sem ser convertida em ausência de falhas.
O comando integrado implementa a seleção dos primeiros cem projetos elegíveis
completos na ordem da busca, que prioriza estrelas. Na execução disponível,
contudo, usamos o piloto já coletado: a coleta adicional foi encerrada por
limite de tempo. Cem candidatos não equivalem a cem incluídos, e a meta de
cem elegíveis completos não foi atingida.

\subsection{Definições operacionais e análise}
Deployment frequency é o número de releases estáveis publicadas dividido
pelas semanas reais da janela, calculadas a partir de sua duração em segundos. A CFR (a) é calculada pela Equação~\ref{eq:cfr}; são falhas as conclusions
\texttt{failure}, \texttt{timed\_out} e \texttt{startup\_failure}, e sucesso a
conclusion \texttt{success}. As demais conclusions são ignoradas. Sem runs
válidos, a CFR fica ausente.
\begin{equation}
\mathrm{CFR}(a)=\frac{\text{falhas}}{\text{falhas}+\text{sucessos}}
\label{eq:cfr}
\end{equation}
Os runs criados na janela são ordenados por \texttt{run\_started\_at} e ID em
cada workflow. A primeira falha após um sucesso observado inicia o episódio;
a duração é \texttt{updated\_at} do próximo sucesso menos
\texttt{run\_started\_at} da primeira falha. Falhas consecutivas mantêm o início.
Sem término antes do fim da janela, o episódio tem censura à direita, e
reportamos sua proporção por repositório. Sequências iniciais sem sucesso
anterior observado têm censura à esquerda separada, sem duração presumida.
Timestamps essenciais inválidos deixam o workflow sem estimativa, com
diagnóstico, evitando unir episódios ao remover somente um run.
A mediana e o intervalo interquartil (IQR), em horas, usam apenas recuperações
observadas, com interpolação linear inclusiva; não são estimativas de sobrevivência.

A coleta de commits compara releases estáveis consecutivas, recuperando a
release anterior à janela quando existir. O protocolo usa a data do autor do
commit. A consulta é paginada e registra comparações inacessíveis ou incompletas;
a primeira release histórica sem anterior é ignorada. O lead time por release
é a publicação menos a data do autor mais antiga entre os commits da comparação.
O valor do repositório é a mediana por release, sem ponderação por número de
commits. Ausência de commits, datas inválidas, comparações incompletas e tempos
negativos não produzem valor; zero observado é válido. Essa variante está
implementada e aguarda validação real. A variante por commit ainda depende de
integração e validação antes da análise conjunta do protocolo.
A classificação usa os cortes fixos do enunciado da disciplina: frequência
Elite a partir de sete releases/semana, High a partir de uma/semana e Medium
a partir de uma/mês; para lead time, os limites são 24, 168 e 720 horas.
CFR usa 15\%, 30\% e 45\%, e recuperação usa 1, 24 e 168 horas.
Os rótulos de classe originalmente previstos nas hipóteses ficam preservados
como registro anterior, mesmo onde divergiam dos cortes agora confirmados.

\subsection{Artefatos e rastreabilidade}
\begingroup\raggedright
O código, as decisões e os testes estão no repositório público:
\url{https://github.com/Victorgabrielcruz/lab-medicao-e-experimentacao-de-software}.\par
O planejamento pode ser consultado no GitHub Projects:
\url{https://github.com/users/Victorgabrielcruz/projects/7}.\par
O relatório técnico está em \texttt{Lab03/docs/relatorio-victor-sprint-01.md}.\par
Dados brutos, cache, credenciais e saídas reais permanecem locais. Hashes da
entrada e contagens agregadas permitem identificar a execução auditada.
\par\endgroup
